\documentclass[11pt]{article}
\usepackage[margin=1in]{geometry}
\usepackage[T1]{fontenc}
\usepackage{lmodern}

\begin{document}

\begin{center}
  {\Large \textbf{Personal Statement}}\\[4pt]
  {\large Ian Chi}\\[2pt]
  Ph.D. Applicant, Department of Electrical Engineering and Computer Sciences\\
  University of California, Berkeley
\end{center}

\vspace{1em}

% Write here.

Most people have their times tables memorized before they turn ten. I learned only a few products, such as the squares, and worked the rest out from them. I still work this way, though not specifically with times tables anymore. I decide what is worth remembering, keep it as a basis from which I can recover the rest when I need it, and rebuild everything else on the fly. Working out a minimal set worth remembering has become one of the ways I come to understand a subject. I find it most interesting when there are incomparable minimal ideas, each enough to recover everything.

That habit is what drew me to mathematics and physics, where a small set of ideas can account for an enormous amount. I majored in both at the University of California, Davis. Computer science came later, through my roommate at the time, who was doing research on automata and regular expressions. He would describe machines that read a string one symbol at a time and can remember only a fixed amount about what they have seen, and the sharp limits that follow on what they can recognize. It was my own question in a new form: what is worth remembering? I started helping him think through the problems he was stuck on, and found the same pleasure I had found in mathematics and physics: a precise question and an argument that settles it.

I came to UC Santa Cruz for a master's degree in computer science to find out whether this new interest would last. In my first quarter I took classes from three theory professors and found that theoretical computer science was the kind of mathematics I had been looking for. I began working with C. Seshadhri on the longest increasing subsequence problem, and what began as an exploration became a real commitment.

The master's degree answered the question I came with, and that is why I am applying now. I have taken a research problem from an outline to a complete proof, and the work has left me with questions larger than a master's degree can hold. I also know what I want to do with them: I plan to stay in academia and become a professor, and a PhD is the next step toward both the research and the teaching that career involves.

Coming to computer science from mathematics and physics, I sometimes encountered advanced ideas before learning the elementary ones beneath them. In my approximation algorithms course, for instance, I learned the expander mixing lemma before learning that the fact that the degrees in a graph sum to twice the number of edges is called the handshake lemma. I often know an idea well before I know its name. This is a consequence of how I learn: I keep a small basis of ideas and rebuild the rest when I need them. Teaching turned out to be how I closed those gaps.

I have tried to contribute to the department as well as take from it. I have tutored undergraduate algorithms and served as a teaching assistant for undergraduate discrete mathematics, where I like to push students into the areas where they are likely unclear, such as empty sets and the difference between ``element of'' and ``subset,'' edge cases that rarely appear in standard exercises. The exchange goes both ways. I bring the habits of mathematical thinking from my first degree to students encountering them for the first time, and their questions send me back to standard topics that I never covered explicitly in a class.

This fall I am hosting the UC Santa Cruz theory seminar, which means finding speakers, introducing them, and making sure the talks are understandable to students who are new to the area. That last part matters to me because I was recently one of those students, often knowing an idea without knowing its name. A talk that assumes the vocabulary can shut out someone who already has the understanding.

Moving between fields has taught me that the questions I care about can survive changes in language and perspective, and that people reach the same mathematics by very different roads. Berkeley's Principles of Community ask for learning opportunities that are open to everyone and for respect for our differences as well as what we share. I have tried to practice both, in the classroom and in the seminar room, and I want to keep doing so at Berkeley and later as a professor: learning from people whose backgrounds differ from mine while helping others find their own ways into mathematics and computer science.

\end{document}
